\chapter[Conclusions \& Future Scope]{Conclusions \& Future Scope}{Conclusions \& Future Scope}\label{res}

\noindent
\rule{0.49\textwidth}{0.75pt} $_{\bigcirc}$ \rule{0.49\textwidth}{0.75pt}\\
\section{Conclusions}
This project developed and evaluated DAVS, a secure, communication-efficient, and audit-ready federated learning framework designed for privacy-sensitive medical imaging environments. The system enables distributed hospital clients to collaboratively train a deep learning model without sharing raw patient data, while addressing the critical challenges of non-IID data heterogeneity, Byzantine poisoning attacks, and the lack of training auditability in conventional federated architectures.

The completed framework integrates the following components into a unified operational pipeline:
\begin{itemize}
    \item Non-IID data partitioning and baseline FedAvg training, establishing that the model converges reliably in a standard federated setting.
    \item Random Projection-based gradient sketching, compressing high-dimensional model updates into compact sketches, significantly reducing per-round bandwidth while preserving similarity structure for trust evaluation.
    \item Data-Aware Verifier Selection (DAVS), an amnesic scoring mechanism that ranks clients using current-round directional alignment (cosine similarity) and magnitude consistency to form a robust validation committee.
    \item PBFT-inspired consensus validation, ensuring that the global model is updated only upon achieving a quorum of validator approvals and rejecting compromised updates.
    \item Permissioned blockchain audit logging, recording per-round scores, committee decisions, vote logs, model hashes, and metrics to provide tamper-proof training provenance.
\end{itemize}

Experimental evaluation under multiple poisoning strategies demonstrated that DAVS sustains accuracy close to the clean baseline while significantly outperforming undefended federated learning under attack, and it simultaneously provides immutable, verifiable provenance of the full training lifecycle.
\noindent
\begin{center}
\rule{0.45\textwidth}{0.75pt} \; $\bigcirc$ \; \rule{0.45\textwidth}{0.75pt}
\end{center}

%\chapter{Conclusions \& Future Scope}



\section{Future Scope}
With the core framework implemented and validated, the following directions represent natural extensions for future research and deployment:
\begin{enumerate}
    \item Extending the framework to larger model architectures (e.g., ResNets and Vision Transformers) and evaluating on diverse clinical datasets to validate generalizability across healthcare scenarios.
    \item Integrating formal differential privacy guarantees into the sketching layer to provide mathematically provable privacy bounds in addition to the obfuscation offered by random projection.
    \item Optimizing the consensus and audit layers for deployment on resource-constrained edge environments and improving throughput under larger client populations.
    \item Investigating robustness under colluding adversaries and more sophisticated adaptive attack strategies, including coordinated poisoning and dynamic participation patterns.
\end{enumerate}


